% A contamination-robust infrared-excess technosignature search around white dwarfs.
% Drafted in the `article` class so it compiles anywhere; the submission version
% targets MNRAS (swap to mnras.cls / \documentclass{mnras}). Numerical results
% marked \TODO are filled in from a science run (`make data && make all`).
\documentclass[11pt]{article}
\usepackage[margin=1in]{geometry}
\usepackage{amsmath,graphicx,natbib,booktabs,hyperref,siunitx}
\newcommand{\TODO}[1]{\textcolor{red}{[#1]}}
\usepackage{xcolor}
\bibliographystyle{plainnat}

\title{A contamination-robust search for infrared-excess technosignatures
around white dwarfs}
\author{[Author list]}
\date{\today}

\begin{document}
\maketitle

\begin{abstract}
Waste heat from a planetary-scale energy-collecting structure (a Dyson sphere or
swarm) is a classic technosignature that would manifest as anomalous infrared
excess. Systematic searches to date have targeted main-sequence stars, most
notably Project Hephaistos, whose candidates proved vulnerable to background
contamination. White dwarfs have been considered only theoretically despite
being uniquely favourable targets: their photospheres are near-blackbody and
fully specified by Gaia astrometry and effective temperature, so an infrared
excess is cleanly defined, and the sole significant natural source of white-dwarf
infrared excess --- a dusty debris disk --- is already catalogued and can be used
as a labelled control. We present the first systematic catalogue search for
infrared-excess technosignatures around white dwarfs, combining the Gaia EDR3
white-dwarf catalogue with CatWISE2020, unWISE, AllWISE and 2MASS photometry. Our
central methodological contribution is a contamination-robust vetting funnel
that adds a Gaia--WISE co-movement test --- requiring the infrared source to
share the white dwarf's large proper motion --- to reject chance-aligned
background interlopers that defeated earlier searches. After subtracting known
debris disks, we report \TODO{$N$} anomalous candidates / a null result, and
place a 95\% occurrence-rate upper limit of \TODO{$f<\ldots$} on white-dwarf
Dyson-swarm-like excess as a function of covering fraction and waste-heat
temperature. We emphasise that broadband photometry cannot uniquely distinguish
a Dyson swarm from warm dust; our candidates are prioritised follow-up targets,
not detections.
\end{abstract}

\section{Introduction}
\label{sec:intro}
% Technosignatures and waste-heat / Dyson structures; the energy argument.
% Project Hephaistos main-sequence search \citep{Suazo2024} and the contamination
% critique \citep{Suazo2024contam}. The white-dwarf case is theoretical only
% \citep{Zackrisson2022,DysonHR2026}. Gap statement: no systematic white-dwarf
% technosignature search exists; debris-disk surveys provide a labelled control.

\section{Data}
\label{sec:data}
% Parent sample: Gaia EDR3 white-dwarf catalogue \citep{GentileFusillo2021},
% Pwd>0.9. Infrared photometry: CatWISE2020 \citep{Marocco2021} primary, unWISE
% \citep{Schlafly2019} independent cross-check, AllWISE quality flags and
% neighbourhood, 2MASS anchor. Control debris-disk catalogues
% \citep{Debes2011,Dennihy2020,Xu2021,MadurgaFavieres2024,MurilloOjeda2026}.
% Volumes and access (TAP/VizieR/CDS X-Match), all catalogue-scale.

\section{Methods}
\label{sec:methods}
\subsection{Photospheric SED prediction}
% Blackbody (self-contained) and Bergeron model-atmosphere predictions of W1/W2,
% anchored on 2MASS J/H/Ks (and Gaia). Robustness via two independent models.
\subsection{Infrared-excess statistic}
The per-band excess significance is
\begin{equation}
  \chi_B = \frac{F_{\mathrm{obs},B} - F_{\mathrm{pred},B}}{\sigma_B},\qquad
  \sigma_B^2 = \sigma_{\mathrm{phot},B}^2 + \sigma_{\mathrm{sys}}^2,
\end{equation}
with a systematic floor $\sigma_{\mathrm{sys}}$ absorbing model and zero-point
error. Selection requires a significant red $W1-W2$ colour excess together with
a significant flux excess in at least one band; the inclusive (OR) band logic
retains cool, $W2$-dominated excesses characteristic of low-temperature
structures.

\subsection{Contamination rejection}
\label{sec:contam}
Our vetting funnel applies, in order: (i) Gaia astrometric quality
(RUWE, parallax signal-to-noise, excess noise); (ii) WISE photometric quality
(contamination/confusion flags, photometric quality, extended-source flag,
saturation, signal-to-noise); (iii) crowding/blending via the AllWISE
neighbourhood multiplicity and unWISE--CatWISE flux agreement; (iv) a
\emph{Gaia--WISE co-movement test}; and (v) extragalactic rejection. The
co-movement test propagates the precise Gaia position to the WISE mean epoch
using the Gaia proper motion and requires the WISE position to follow, and the
CatWISE2020 proper motion (where measured) to agree:
\begin{equation}
  \Delta\theta = \left| \vec{x}_{\rm WISE} - \left[\vec{x}_{\rm Gaia} +
  \vec{\mu}_{\rm Gaia}\,(t_{\rm WISE}-t_{\rm Gaia})\right] \right| < \Delta\theta_{\max}.
\end{equation}
Because white dwarfs are nearby and fast-moving, a static or differently-moving
infrared source at the catalogued position is revealed as a chance-aligned
background object --- the dominant failure mode of earlier broadband searches.

\subsection{Discriminating technosignatures from debris disks}
We fit a single-temperature blackbody $(T_{\rm dust},\tau)$ to each excess,
calibrate the empirical white-dwarf debris-disk locus from the control
catalogues, subtract known disks, and flag locus outliers (too cool, or
$\tau$ approaching unity) as candidates. We state explicitly that photometry
alone is degenerate between warm dust and a partial swarm; degeneracy-breaking
requires optical variability (ZTF) or full-SED follow-up.

\section{Results}
\label{sec:results}
% Contamination funnel survivor counts (Fig.~\ref{fig:funnel}); recovery of the
% published debris-disk samples as a pipeline validation; candidate table or
% null result; occurrence-rate upper limit and completeness (Fig.~\ref{fig:cc}).

\begin{figure}[t]
  \centering
  % \includegraphics[width=0.7\linewidth]{../results/figures/funnel.pdf}
  \caption{Contamination funnel: white dwarfs surviving each successive cut.}
  \label{fig:funnel}
\end{figure}

\begin{figure}[t]
  \centering
  % \includegraphics[width=0.7\linewidth]{../results/figures/color_color_locus.pdf}
  \caption{Excess $(T_{\rm dust},\tau)$ with the empirical debris-disk locus;
  anomaly candidates lie outside it.}
  \label{fig:cc}
\end{figure}

\section{Discussion}
\label{sec:discussion}
% Dust/swarm degeneracy and what variability/SED-shape adds; comparison to the
% Zackrisson (2022) detectability predictions and to Hephaistos main-sequence
% limits; prospects with JWST/Rubin.

\section{Conclusions}
\label{sec:conclusions}

\section*{Data and code availability}
All code is openly available (Zenodo DOI \TODO{}); the analysis reproduces
end-to-end from public catalogues via the committed pipeline.

\bibliography{refs}
\end{document}
